\section*{Chapitre \ref{ch:b3:conformal} --- Applications
conformes et théorème de l'application de Riemann}

\begin{solution}{exo:b3:conformal:1}
(a) $\varphi$ et $\psi(w) = \iu\frac{1+w}{1-w}$ se composent
en l'identité dans les deux ordres (calcul direct)~;
$\varphi$ envoie $\mathbb H$ dans $\mathbb D$ et $\psi$ en
retour (l'\cref{ex:b3:conformal:mobius}). Valeurs~:
$\varphi(\iu) = 0$, $\varphi(0) = -1$, $\varphi(1) =
\frac{1 - \iu}{1 + \iu} = -\iu$, et $\varphi(z) \to 1$ quand
$z \to \infty$.

(b) Les cercles et droites sont les ensembles de solutions de
$\alpha\abs z^2 + \bar\beta z + \beta\bar z + \gamma = 0$
($\alpha, \gamma \in \R$, $\beta \in \C$, $\abs\beta^2 >
\alpha\gamma$)~: $\alpha \neq 0$ cercles, $\alpha = 0$
droites. En substituant $z = 1/w$ et multipliant par $\abs
w^2$~: $\gamma\abs w^2 + \beta w + \bar\beta\bar w +
\alpha = 0$ --- même famille. Les applications affines
préservent clairement la famille, et toute application de
Möbius est une composition d'applications affines et d'une
inversion ($\frac{az+b}{cz+d} = \frac ac + \frac{bc -
ad}{c}\cdot\frac1{cz + d}$ pour $c \neq 0$).
\end{solution}

\begin{solution}{exo:b3:conformal:2}
(a) Schwarz donne $\abs{f(a)} \leq \abs a$ avec égalité (les
deux côtés $= \abs a$)~: le cas d'égalité force $f(z) =
\eu^{\iu\theta}z$, et $f(a) = a$ fixe $\eu^{\iu\theta} = 1$.

(b) Soient $a \neq b$ des points fixes et $g =
\varphi_a\circ f\circ\varphi_{-a} \in
\operatorname{Aut}(\mathbb D)$ --- en utilisant
le~\cref{thm:b3:conformal:autdisc} pour $\varphi_{\pm a}$.
Alors $g(0) = \varphi_a(f(a)) = 0$ et $g(c) = c$ pour $c =
\varphi_a(b) \neq 0$~: par (a), $g = \mathrm{id}$, donc $f =
\varphi_{-a}\circ\varphi_a = \mathrm{id}$.
\end{solution}

\begin{solution}{exo:b3:conformal:3}
Fixer $z$ et poser $g = \varphi_{f(z)}\circ
f\circ\varphi_{-z}$~: \omterm{def:b3:holomorphic:holo}{holomorphe} $\mathbb D \to \mathbb D$
avec $g(0) = 0$, donc $\abs{g'(0)} \leq 1$ (Schwarz). Règle
de la chaîne avec $\varphi_a'(\zeta) = \frac{1 - \abs
a^2}{(1 - \bar a\zeta)^2}$~:
\[
g'(0) = \varphi_{f(z)}'\bigl(f(z)\bigr)\cdot f'(z)\cdot
\varphi_{-z}'(0)
= \frac{1}{1 - \abs{f(z)}^2}\cdot f'(z)\cdot(1 - \abs z^2),
\]
d'où l'inégalité de Schwarz--Pick. L'égalité en un $z$ rend
$g$ une rotation, d'où $f =
\varphi_{-f(z)}\circ(\text{rotation})\circ\varphi_z \in
\operatorname{Aut}(\mathbb D)$ --- et alors l'égalité vaut
partout (calculer, ou réappliquer avec les rôles de $f,
f^{-1}$ échangés). Les auto-applications \omterm{def:b3:holomorphic:holo}{holomorphes} du
disque sont $1$-lipschitziennes pour la métrique
hyperbolique $\frac{2\abs{\dd z}}{1 - \abs z^2}$~; les
automorphismes en sont les isométries.
\end{solution}

\begin{solution}{exo:b3:conformal:4}
(a) $z \mapsto \eu^z$~: envoie $\{0 < \operatorname{Im}z <
\pi\}$ bijectivement sur $\mathbb H$ ($\eu^{x+\iu y} =
\eu^x\eu^{\iu y}$~: \omterm{def:b3:modules:free}{module libre}, argument $y \in
\intoo0\pi$), \omterm{def:b3:holomorphic:holo}{holomorphe} de dérivée non nulle et inverse
\omterm{def:b3:holomorphic:holo}{holomorphe} ($\log$ principal).
(b) $z \mapsto z^2$ double les arguments~: le quadrant
\omterm{def:b3:topology:topology}{ouvert} $\{0 < \arg z < \frac\pi2\}$ s'envoie conformément
sur $\mathbb H$ (inverse~: racine carrée principale).
(c) $z \mapsto \frac{1 + z}{1 - z}$ envoie $\mathbb D$ sur
le demi-plan droit et préserve la symétrie haut/bas~: elle
envoie le demi-disque supérieur sur le premier quadrant~;
puis élever au carré, par (b), pour atteindre $\mathbb H$~:
$z \mapsto \bigl(\frac{1 + z}{1 - z}\bigr)^2$.
(d) Le \omterm{thm:b3:conformal:autdisc}{facteur de Blaschke} $\varphi_{1/2}(z) = \frac{z -
\frac12}{1 - \frac z2}$~: $\varphi_{1/2}(\tfrac12) = 0$ et
$\varphi_{1/2}'(\tfrac12) = \frac{1 - \frac14}{(1 -
\frac14)^2} = \frac43 > 0$.
\end{solution}

\begin{solution}{exo:b3:conformal:5}
(a) Une \omterm{def:b3:conformal:conformal}{application conforme} $\C \to \mathbb D$ (ou $\mathbb
H$, après composition avec Cayley) est une fonction entière
bornée~: constante par Liouville --- pas bijective. Une
conforme $\mathbb D \to \C$ aurait un inverse conforme $\C
\to \mathbb D$~: même contradiction.

(b) Le carré est convexe, d'où \omterm{def:b3:conformal:simplyconnected}{simplement connexe}, et
propre~: conformément $\mathbb D$
(le~\cref{thm:b3:conformal:rmt}). Le plan découpé
$\C\setminus\intoc{-\infty}0$ est étoilé par rapport à $1$
(les segments depuis $1$ évitent la coupure), d'où
\omterm{def:b3:conformal:simplyconnected}{simplement connexe}, et propre~: conformément $\mathbb D$.
Le disque percé~: si $g \colon \mathbb D\setminus\{0\} \to
\mathbb D$ était conforme, $g$ est bornée, donc $0$ est
éliminable (le~\cref{thm:b3:residues:riemanncw})~: $g$
s'étend en $\tilde g \colon \mathbb D \to \mathbb D$, et
$\tilde g(0)$, étant dans l'image ouverte $g(\mathbb
D\setminus\{0\}) = \mathbb D$, est aussi $g(w)$ pour un $w
\neq 0$~; deux \omterm{def:b3:topology:topology}{voisinages} disjoints de $0$ et $w$ ont des
images \omterm{def:b3:topology:topology}{ouvertes} partageant la valeur $\tilde g(0)$, d'où
partageant \emph{d'autres} valeurs aussi (\omterm{def:b3:topology:topology}{ouverts})~: $g$
prend une certaine valeur deux fois sur $\mathbb
D\setminus\{0\}$ --- contredisant l'injectivité. L'anneau
$A = \{1 < \abs z < 2\}$~: supposons $F \colon \mathbb D
\to A$ conforme. $F$ est sans zéro sur le \omterm{def:b3:conformal:simplyconnected}{simplement
connexe} $\mathbb D$, donc $F = \eu^L$ pour $L$ \omterm{def:b3:holomorphic:holo}{holomorphe}
(la~\cref{def:b3:conformal:simplyconnected}). Soit $\sigma$
le cercle $\abs z = \frac32$ dans $A$ et $\gamma =
F^{-1}\circ\sigma$, un chemin fermé dans $\mathbb D$~; alors
\[
1 = \operatorname{Ind}_\sigma(0)
= \frac1{2\iu\pi}\int_{F\circ\gamma}\frac{\dd w}{w}
= \frac1{2\iu\pi}\int_\gamma\frac{F'}{F}
= \frac1{2\iu\pi}\int_\gamma L' = 0
\]
($L'$ a une primitive)~: contradiction. Ni le disque percé
ni l'anneau n'est un disque déguisé.
\end{solution}

\begin{solution}{exo:b3:conformal:6}
(a) $T(w) = \frac{w - 1}{w + 1}$ envoie $\{\operatorname{Re}w
> 0\}$ conformément sur $\mathbb D$ (Cayley tourné~: $\abs{w
- 1} < \abs{w + 1}$ ssi $\operatorname{Re}w > 0$), avec
$T(1) = 0$. Pour $f \in \mathcal F$, $g = T\circ f\colon
\mathbb D \to \mathbb D$ est \omterm{def:b3:holomorphic:holo}{holomorphe} avec $g(0) = 0$~:
Schwarz donne $\abs{g(z)} \leq \abs z$.

(b) En invertissant $T$~: $f = \frac{1 + g}{1 - g}$, donc
\[
\abs{f(z)} \leq \frac{1 + \abs{g(z)}}{1 - \abs{g(z)}}
\leq \frac{1 + \abs z}{1 - \abs z} :
\]
localement bornée (uniformément sur $\abs z \leq r < 1$).
L'égalité en $z_0 \neq 0$ force $\abs{g(z_0)} = \abs{z_0}$
et l'alignement~: $g$ une rotation, c'est-à-dire $f(z) =
\frac{1 + \eu^{\iu\theta}z}{1 - \eu^{\iu\theta}z}$ --- les
\emph{extrémales de Herglotz}, \omterm{def:b3:conformal:conformal}{applications conformes} sur le
demi-plan droit.
\end{solution}

\begin{solution}{exo:b3:conformal:7}
(i) La \omterm{def:b3:groups:simple}{simple} connexité entre exactement deux fois, via
l'existence de racines carrées \omterm{def:b3:holomorphic:holo}{holomorphes} de fonctions sans
zéro (la~\cref{def:b3:conformal:simplyconnected})~: à
l'étape~0 (la racine de $z - b$) et à l'étape~2 (la racine
de $\varphi_a\circ f$).
(ii) $\Omega \neq \C$ fournit le point $b$ de l'étape~0 ---
sans lui la famille $\mathcal F$ est vide d'applications
injectives bornées (Liouville~: toute $f \colon \C \to
\mathbb D$ \omterm{def:b3:holomorphic:holo}{holomorphe} est constante). (iii) La connexité est
utilisée chaque fois que le théorème d'identité ou Hurwitz
(l'\cref{exo:b3:residues:8}) parle~: la limite extrémale est
«~injective ou constante~», et la constance est exclue par
$s > 0$~; aussi dans «~dérivée nulle implique constante~».
Échec pour $\C$~: étape~0 impossible, et la conclusion est
fausse (l'\cref{exo:b3:conformal:5}(a)). Échec pour
l'anneau~: pas \omterm{def:b3:conformal:simplyconnected}{simplement connexe} --- la construction par
racine carrée casse (p.\ ex.\ $z$ elle-même, sans zéro sur
$A$, n'a pas de racine carrée \omterm{def:b3:holomorphic:holo}{holomorphe}~: le même calcul
d'\omterm{def:b3:holomorphic:index}{indice} que dans l'\cref{exo:b3:conformal:5}(b) avec
$\frac12\int_\sigma\frac{\dd z}z \notin 2\iu\pi\Z$) --- et
la conclusion est fausse aussi.
\end{solution}

\begin{solution}{exo:b3:conformal:8}
En substituant le développement de Fourier de $g$ dans
l'intégrale de Poisson et en utilisant
$\frac1{2\pi}\int P_r(\varphi - t)\eu^{\iu nt}\dd t =
r^{\abs n}\eu^{\iu n\varphi}$ (lu sur la série de $P_r$)~:
$u(r\eu^{\iu\varphi}) = \sum_nc_n(g)\,r^{\abs
n}\eu^{\iu n\varphi}$, l'interversion justifiée par
convergence normale ($\abs{c_n} \leq \norm g_\infty$, $r <
1$).

(a) $g = \cos t$~: $c_{\pm1} = \frac12$, donc $u =
r\cos\varphi = \operatorname{Re}z = x$ --- bien harmonique
avec les bonnes valeurs au bord.

(b) $\cos^2t = \frac12 + \frac{\cos 2t}2$~: $u = \frac12 +
\frac{r^2\cos2\varphi}2 = \frac12 +
\frac{\operatorname{Re}(z^2)}2 = \frac12 + \frac{x^2 -
y^2}{2}$.

(c) $g = \mathbf 1_{(0,\pi)}$ (demi-cercle supérieur)~: $c_0
= \frac12$ et $c_n = \frac{1 - (-1)^n}{2\iu\pi n}$ pour $n
\neq 0$, donc
\[
u(r\eu^{\iu\varphi}) = \frac12 + \frac2\pi\sum_{k\geq0}
\frac{r^{2k+1}\sin\bigl((2k+1)\varphi\bigr)}{2k + 1},
\qquad u(0) = \frac12 :
\]
le \omterm{ex:b3:groups:actions}{centre} voit exactement la moyenne des données au bord ---
la propriété de la moyenne en personne.
\end{solution}

\begin{solution}{exo:b3:conformal:9}
De $(1-r)^2 \leq 1 - 2r\cos\theta + r^2 \leq (1+r)^2$~:
\[
\frac{1-r}{1+r} = \frac{1 - r^2}{(1+r)^2} \leq P_r(\theta)
\leq \frac{1 - r^2}{(1 - r)^2} = \frac{1+r}{1-r} .
\]
Pour $u \geq 0$ harmonique sur $\mathbb D$ et $s < 1$~:
$u_s(z) = u(sz)$ est harmonique sur un \omterm{def:b3:topology:topology}{voisinage} de
$\bar{\mathbb D}$, d'où égale son intégrale de Poisson
(le~\cref{thm:b3:conformal:poisson}, unicité, appliqué à ses
propres valeurs au bord)~; en encadrant le noyau et en
utilisant la moyenne $\frac1{2\pi}\int u_s(\eu^{\iu t})\dd t
= u(0)$~:
\[
\frac{1-r}{1+r}\,u(0) \leq u(s\,r\eu^{\iu\varphi}) \leq
\frac{1+r}{1-r}\,u(0).
\]
Faire $s \to 1^-$ à $r\eu^{\iu\varphi}$ fixe (continuité de
$u$)~: les inégalités de Harnack. Si $u$ est harmonique sur
$\C$ avec $u \geq m$~: appliquer Harnack à $u - m$ sur les
disques $D(0, R)$, c'est-à-dire à $z \mapsto u(Rz) - m$~:
pour $z$ fixe et $r = \abs z/R \to 0$, les deux bornes
tendent vers $u(0) - m$~: $u(z) = u(0)$ --- constante (un
Liouville bilatère à partir d'une borne unilatère).
\end{solution}

\begin{solution}{exo:b3:conformal:10}
Localement, $u = \operatorname{Re}F$ avec $F$ \omterm{def:b3:holomorphic:holo}{holomorphe}
(la~\cref{prop:b3:conformal:harmonicholo}), donc $u\circ f =
\operatorname{Re}(F\circ f)$ est harmonique là où elle est
définie~: l'harmonicité est conformément invariante. Sur
$\mathbb H$~: le logarithme principal donne $\log z =
\ln\abs z + \iu\arg z$ \omterm{def:b3:holomorphic:holo}{holomorphe} sur $\mathbb H$, donc $u =
\frac1\pi\arg z = \operatorname{Im}\bigl(\frac1\pi\log
z\bigr)$ est harmonique, avec limites au bord~: pour $x >
0$, $\arg \to 0$, $u \to 0$~; pour $x < 0$, $\arg \to \pi$,
$u \to 1$~: les données $\mathbf 1_{\intoo{-\infty}0}$ en
tout $x \neq 0$. En transportant par l'\omterm{ex:b3:conformal:mobius}{application de
Cayley} (qui envoie $\mathbb D \to \mathbb H$ après
inversion et fait correspondre le demi-cercle supérieur à
l'axe négatif, à la rotation près fixée en suivant trois
points du bord), $u\circ(\text{Cayley})$ résout le problème
du disque de l'\cref{exo:b3:conformal:8}(c)~; l'évaluation
au \omterm{ex:b3:groups:actions}{centre} retrouve $u = \frac12$ là, et la forme fermée
$\frac1\pi\arg$ peut être contrôlée contre la série en
sommant $\sum\frac{r^{2k+1}\sin((2k+1)\varphi)}{2k+1} =
\frac12\arctan\frac{2r\sin\varphi}{1 - r^2}$-type
identités --- la route élémentaire vers la même réponse.
\end{solution}

\begin{solution}{exo:b3:conformal:11}
(a) Soient $a \neq b$ fixes. En conjuguant par $\varphi_a(z)
= \frac{z - a}{1 - \bar az}$ (un automorphisme échangeant
$a$ et $0$), $g = \varphi_a\circ f\circ\varphi_a^{-1}$ fixe
$0$ et le point $c = \varphi_a(b) \neq 0$. Schwarz~:
$\abs{g(z)} \leq \abs z$, et en $z = c$ l'égalité vaut
($g(c) = c$)~: le cas d'égalité force $g(z) = \lambda z$
avec $\abs\lambda = 1$, et $\lambda c = c$ donne $\lambda =
1$~: $g = \mathrm{id}$, d'où $f = \mathrm{id}$.

(b) $h(z) = f(z)/z$ (\omterm{thm:b3:residues:riemanncw}{singularité éliminable} en $0$) est
\omterm{def:b3:holomorphic:holo}{holomorphe} sur $\mathbb D$ avec $\abs h \leq 1$
(Schwarz)~; $\abs h < 1$ partout, sinon le principe du
maximum (maximum intérieur de $\abs h$) rendrait $h$ une
constante unimodulaire, c'est-à-dire $f$ une rotation ---
exclu. Sur le compact $\bar D(0,r)$, $c_r = \max\abs h <
1$, donc $\abs{f(z)} \leq c_r\abs z$ là~; de plus $f$ envoie
$\bar D(0,r)$ dans lui-même ($c_r\abs z \leq r$), donc la
borne itère~: $\abs{f^{\circ n}(z)} \leq c_r^n\,r \to 0$
uniformément sur $\bar D(0, r)$.

(c) Points fixes de $\frac{z^2 + z}2$~: $z^2 + z = 2z$ ssi
$z(z - 1) = 0$~; seul $z = 0$ est dans $\mathbb D$ ($z = 1$
est sur le bord). Pas une rotation ($f'(0) = \frac12$),
donc les \omterm{def:b3:groups:action}{orbites} tendent vers $0$~; quantitativement $f(z)
= \frac z2(1 + z)$ donne $\abs{f(z)} \leq \frac{3}{4}\abs
z$ sur $\abs z \leq \frac12$, et une fois l'\omterm{def:b3:groups:action}{orbite} petite,
$\abs{f(z)} \approx \frac{\abs z}2$~: asymptotiquement
géométrique de raison $f'(0) = \frac12$. Depuis $z_0 =
\frac12$~: $z_1 = \frac38$, $z_2 \approx 0{,}258$, $z_3
\approx 0{,}162$ --- une demi-vie par pas, comme prédit par
le multiplicateur.
\end{solution}

\begin{solution}{exo:b3:conformal:12}
(a) $\Delta u = 6x - 6x + 0 = 0$. Cauchy--Riemann exige $v_y
= u_x = 3x^2 - 3y^2$ et $v_x = -u_y = 6xy - 2$. En intégrant
la première en $y$~: $v = 3x^2y - y^3 + c(x)$~; en
branchant dans la seconde~: $6xy + c'(x) = 6xy - 2$, donc
$c(x) = -2x + C$. Ainsi $v = 3x^2y - y^3 - 2x + C$ et
\[
f = u + \iu v = (x^3 - 3xy^2) + \iu(3x^2y - y^3)
+ 2y - 2\iu x + \iu C = z^3 - 2\iu z + \iu C .
\]

(b) La forme $\omega = -u_y\,\dd x + u_x\,\dd y$ est fermée
précisément parce que $\Delta u = 0$ ($\partial_y(-u_y) =
-u_{yy} = u_{xx} = \partial_x(u_x)$). Sur un \omterm{def:b3:topology:topology}{ouvert} étoilé
le lemme de Poincaré (le~\cref{thm:b3:forms:poincare}~; ou
la construction de primitive
du~\cref{thm:b3:holomorphic:cauchy} appliquée à
l'\omterm{def:b3:holomorphic:holo}{holomorphe} $u_x - \iu u_y$) fournit $v$ avec $\dd v =
\omega$, c'est-à-dire le système de Cauchy--Riemann~: $u +
\iu v$ est \omterm{def:b3:holomorphic:holo}{holomorphe}. Deux conjuguées diffèrent d'une
fonction de gradient nul~: une constante (connexité).

(c) Pour $u = \ln\abs z$ sur $\C^*$~: $\omega = -u_y\dd x +
u_x\dd y = \frac{-y\,\dd x + x\,\dd y}{x^2 + y^2} =
\omega_\theta$, la forme angulaire
(l'\cref{ex:b3:forms:angular}), dont l'intégrale le long du
cercle unité est $2\pi \neq 0$~: non exacte, donc aucune
conjuguée globale n'existe --- une conjuguée serait une
détermination \omterm{def:b3:topology:continuity}{continue} de l'argument, et le \omterm{def:b3:holomorphic:index}{nombre de tours}
est exactement l'obstruction. Localement (sur tout
sous-domaine étoilé), $v = \arg z$ marche et $u + \iu v =
\log z$~: l'échec est global, non local.
\end{solution}

\begin{solution}{pb:b3:conformal:1}
\textbf{1.} $g$ est injective et \omterm{def:b3:holomorphic:holo}{holomorphe}~; sur $C_\rho$
($\rho > 1$) elle est lisse, et l'aire encerclée est
$\frac1{2\iu}\oint_{g(C_\rho)}\bar\zeta\,\dd\zeta$ (la
formule d'aire de Green--Riemann de deuxième année,
appliquée avec orientation positive). En substituant $\zeta
= g(w)$, $w = \rho\eu^{\iu\theta}$~:
\[
A_\rho = \frac1{2\iu}\int_{C_\rho}\overline{g(w)}\,g'(w)\dd
w = \frac\rho2\int_0^{2\pi}\overline{g(\rho\eu^{\iu\theta})}
\,g'(\rho\eu^{\iu\theta})\,\eu^{\iu\theta}\,\dd\theta .
\]
Insérer $\bar g = \rho\eu^{-\iu\theta} + \bar b_0 +
\sum_m\bar b_m\rho^{-m}\eu^{\iu m\theta}$ et $g' = 1 -
\sum_nnb_n\rho^{-n-1}\eu^{-\iu(n+1)\theta}$~: après
multiplication par $\eu^{\iu\theta}$, seuls les produits de
fréquence nulle survivent à l'intégration en $\theta$ (la
convergence normale justifie le travail terme à terme)~: la
paire $(\rho\eu^{-\iu\theta})\cdot1$ contribue $2\pi\rho$,
et chaque paire $\bar b_n\rho^{-n}\eu^{\iu
n\theta}\cdot(-nb_n\rho^{-n-1}\eu^{-\iu(n+1)\theta})$
contribue $-2\pi n\abs{b_n}^2\rho^{-2n-1}$~:
\[
A_\rho = \pi\rho^2 - \pi\sum_{n\geq1}n\abs{b_n}^2\rho^{-2n}.
\]

\textbf{2.} Les aires sont non négatives~: $\sum_{n\leq
N}n\abs{b_n}^2\rho^{-2n} \leq \rho^2$ pour tout $N$~; faire
$\rho \to 1^+$ puis $N \to \infty$~:
$\sum_{n\geq1}n\abs{b_n}^2 \leq 1$. L'égalité dans
$\abs{b_1} \leq 1$ force tous les autres $b_n = 0$~: $g(w) =
w + b_0 + \eu^{\iu\alpha}/w$, qui s'envoie sur le
complémentaire d'un segment de longueur $4$ (application de
type Joukowski)~: les extrémales.

\textbf{3.} $f(z)/z = 1 + a_2z + \cdots$ est \omterm{def:b3:holomorphic:holo}{holomorphe} et
sans zéro sur $\mathbb D$ ($f$ ne s'annule qu'en $0$,
simplement~: injectivité), d'où aussi $f(z^2)/z^2$, qui a
une racine carrée \omterm{def:b3:holomorphic:holo}{holomorphe} $\varphi$ avec $\varphi(0) =
1$ (la~\cref{def:b3:conformal:simplyconnected}~; $\mathbb D$
est convexe). Alors $h(z) = z\varphi(z^2)$ satisfait $h^2 =
f(z^2)$, $h(z) = z\bigl(1 + \frac{a_2}2z^2 + \cdots\bigr)$
(série binomiale pour la racine), et $h(z) = z\varphi(z^2)$
est impaire par construction ($\varphi(z^2)$ est paire).
Injectivité~: $h(z) = h(z')$ donne $f(z^2) = f(z'^2)$, donc
$z^2 = z'^2$, c'est-à-dire $z' = \pm z$~; si $z' = -z$ alors
l'imparité donne $h(z) = -h(z)$, donc $h(z) = 0$, forçant
$z = 0$ ($\varphi$ sans zéro)~: $z = z' = 0$.

\textbf{4.} $g(w) = 1/h(1/w)$~: pour $\abs w > 1$, $1/w \in
\mathbb D\setminus\{0\}$ et $h \neq 0$ là~: bien définie,
injective (composition d'injections), avec développement
\[
g(w) = \frac{1}{\frac1w\bigl(1 + \frac{a_2}{2w^2} +
\cdots\bigr)} = w\Bigl(1 - \frac{a_2}{2w^2} + \cdots\Bigr)
= w - \frac{a_2}{2}\,w^{-1} + \cdots :
\]
de la forme de la partie~I avec $b_1 = -\frac{a_2}2$. Le
théorème de l'aire donne $\abs{\frac{a_2}2} \leq 1$~:
$\abs{a_2} \leq 2$.

\textbf{5.} $k(z) = \frac z{(1-z)^2} = z\sum_{m\geq0}(m +
1)z^m = \sum_{n\geq1}nz^n$~: coefficients $a_n = n$, donc
$a_2 = 2$. Univalence et image~: $k = \frac14\bigl[w^2 -
1\bigr]$ avec $w = \frac{1 + z}{1 - z}$, \omterm{def:b3:conformal:conformal}{application
conforme} de $\mathbb D$ sur le demi-plan droit~; $w^2$
envoie ce demi-plan conformément sur
$\C\setminus\intoc{-\infty}0$~; puis $\frac{\cdot - 1}4$
donne $\C\setminus\intoc{-\infty}{-\frac14}$~: injective à
chaque étape, image comme annoncé.

\textbf{6.} $F = \frac{cf}{c - f}$~: comme $c \notin
f(\mathbb D)$, le dénominateur ne s'annule jamais~: $F$ est
\omterm{def:b3:holomorphic:holo}{holomorphe}, et injective ($w \mapsto \frac{cw}{c - w}$ est
Möbius, injective hors de $w = c$). Développement~: avec $f
= z + a_2z^2 + \cdots$,
\[
F = f\cdot\frac1{1 - f/c} = \bigl(z + a_2z^2\bigr)\Bigl(1 +
\frac zc\Bigr) + O(z^3)
= z + \Bigl(a_2 + \frac1c\Bigr)z^2 + O(z^3):
\]
$F \in \mathcal S$ avec $A_2 = a_2 + \frac1c$.

\textbf{7.} Bieberbach deux fois~: $\abs{a_2} \leq 2$ et
$\abs{a_2 + \frac1c} \leq 2$, donc $\abs{\frac1c} \leq 4$~:
$\abs c \geq \frac14$. Toute valeur omise est hors de
$D(0,\frac14)$, c'est-à-dire $f(\mathbb D) \supseteq D(0,
\frac14)$. Netteté~: la fonction de Koebe omet $-\frac14$
(question~5).

\textbf{8.} Soit $d = d(z_0, \partial\Omega)$ et $\psi =
\varphi^{-1} \colon \mathbb D \to \Omega$, $\psi(0) = z_0$,
$\psi'(0) = 1/\varphi'(z_0) > 0$. La normalisation $\tilde
f(z) = \frac{\psi(z) - z_0}{\psi'(0)}$ est dans $\mathcal
S$, donc son image contient $D(0, \frac14)$~; en
re-échelonnant, $\Omega = \psi(\mathbb D) \supseteq
D\bigl(z_0, \tfrac{\psi'(0)}4\bigr)$, donc $d \geq
\frac{\psi'(0)}4 = \frac1{4\varphi'(z_0)}$~: l'inégalité de
gauche. Pour la droite~: $\varphi\circ(z_0 + d\,\cdot)$
envoie $\mathbb D$ dans $\mathbb D$ avec $0 \mapsto 0$~;
Schwarz borne sa dérivée en $0$~: $d\,\varphi'(z_0) \leq 1$,
c'est-à-dire $\frac1{\varphi'(z_0)} \geq d$ --- ensemble
\[
d \leq \frac1{\varphi'(z_0)} \leq 4d .
\]
(L'inégalité de gauche comme affichée dans l'énoncé est la
même chaîne réarrangée.)

\textbf{9.} Pour la fonction de Koebe $a_n = n$~: égalité
dans toute la conjecture~; ses rotations
$\eu^{-\iu\theta}k(\eu^{\iu\theta}z) = \sum
n\eu^{\iu(n-1)\theta}z^n$ ont aussi $\abs{a_n} = n$. En
poussant la question~4~: $h = z + \frac{a_2}2z^3 + c_5z^5 +
\cdots$ avec $h^2 = f(z^2)$~; en comparant les coefficients
de $z^6$~: $2c_5 + \frac{a_2^2}4 = a_3$, donc $c_5 =
\frac{a_3}2 - \frac{a_2^2}8$. Alors
\[
g(w) = \frac1{h(1/w)} = w\Bigl(1 - \frac{a_2}{2w^2} +
\Bigl(\frac{a_2^2}4 - c_5\Bigr)\frac1{w^4} + \cdots\Bigr)
= w - \frac{a_2}2w^{-1} + \Bigl(\frac{3a_2^2}8 -
\frac{a_3}2\Bigr)w^{-3} + \cdots,
\]
et le théorème de l'aire ($\sum n\abs{b_n}^2 \leq 1$) donne
\[
\Bigl|\frac{a_2}2\Bigr|^2 + 3\,\Bigl|\frac{3a_2^2}8 -
\frac{a_3}2\Bigr|^2 \leq 1 .
\]
Contrôle Koebe ($a_2 = 2$, $a_3 = 3$)~: $\abs{\frac{a_2}2}^2
= 1$ et $\frac{3\cdot4}8 - \frac32 = 0$~: total exactement
$1$ --- extrémale, comme il se doit.

\textbf{10.} $\varphi$ est un automorphisme du disque avec
$\varphi(0) = z_0$, donc $f\circ\varphi$ est univalente sur
$\mathbb D$ (composition d'injections), et $f'(z_0) \neq 0$
(la~\cref{def:b3:conformal:conformal})~: $F$ est bien
définie et univalente. $F(0) = 0$. Règle du quotient~:
\[
\varphi'(z) = \frac{(1 + \bar z_0z) - (z + z_0)\bar z_0}
{(1 + \bar z_0z)^2}
= \frac{1 - \abs{z_0}^2}{(1 + \bar z_0z)^2},
\qquad \varphi'(0) = 1 - \abs{z_0}^2 ,
\]
donc $(f\circ\varphi)'(0) = f'(z_0)(1 - \abs{z_0}^2)$ et
$F'(0) = 1$~: $F \in \mathcal S$.

\textbf{11.} Avec $G = f\circ\varphi$~: $G'' =
f''(\varphi)\,\varphi'^2 + f'(\varphi)\,\varphi''$, et
$\varphi''(z) = -2\bar z_0(1 - \abs{z_0}^2)(1 + \bar
z_0z)^{-3}$ donne $\varphi''(0) = -2\bar z_0(1 -
\abs{z_0}^2)$. D'où
\[
A_2 = \frac{G''(0)}{2f'(z_0)(1 - \abs{z_0}^2)}
= \frac12\Bigl[\bigl(1 - \abs{z_0}^2\bigr)
\frac{f''(z_0)}{f'(z_0)} - 2\bar z_0\Bigr] .
\]
Bieberbach pour $F$ (question~4)~: $\abs{A_2} \leq 2$,
c'est-à-dire $\bigl|(1 - \abs{z_0}^2)\frac{f''(z_0)}{f'(z_0)}
- 2\bar z_0\bigr| \leq 4$. Multiplier par $z_0/(1 -
\abs{z_0}^2)$, de \omterm{def:b3:modules:module}{module} $r/(1 - r^2)$ pour $z_0 =
r\eu^{\iu\theta}$, et utiliser $z_0\bar z_0 = r^2$~:
\[
\Bigl|\,z_0\frac{f''(z_0)}{f'(z_0)} -
\frac{2r^2}{1 - r^2}\Bigr| \leq \frac{4r}{1 - r^2} .
\]

\textbf{12.} Un nombre complexe à distance $\rho$ du point
réel $c$ a une partie réelle dans $\intcc{c - \rho}{c +
\rho}$~: appliquer à $c = \frac{2r^2}{1-r^2}$, $\rho =
\frac{4r}{1-r^2}$.

\textbf{13.} Pour $g$ $\mathcal C^1$ non nulle~:
$\log\abs g = \frac12\log(g\bar g)$, donc $\frac{\dd}{\dd
t}\log\abs g = \frac{g'\bar g + g\bar g'}{2\abs g^2} =
\operatorname{Re}\frac{g'}g$. Avec $g(t) =
f'(t\eu^{\iu\theta})$ (sans zéro~: univalence), $g'(t) =
\eu^{\iu\theta}f''(t\eu^{\iu\theta})$, donc en $z =
t\eu^{\iu\theta}$~:
\[
\frac{\dd}{\dd t}\log\bigl|f'(t\eu^{\iu\theta})\bigr|
= \operatorname{Re}\Bigl(\eu^{\iu\theta}
\frac{f''}{f'}\Bigr)
= \frac1t\operatorname{Re}\Bigl(z\frac{f''}{f'}\Bigr)
\in \Bigl[\frac{2t - 4}{1 - t^2},\
\frac{2t + 4}{1 - t^2}\Bigr]
\]
par la question~12 au rayon $t$. Comme $\frac{\dd}{\dd
t}\log\frac{1+t}{(1-t)^3} = \frac1{1+t} + \frac3{1-t} =
\frac{2t+4}{1-t^2}$ et $\frac{\dd}{\dd
t}\log\frac{1-t}{(1+t)^3} = \frac{-1}{1-t} - \frac3{1+t} =
\frac{2t-4}{1-t^2}$, en intégrant de $0$ à $r$ (les trois
fonctions s'annulent en $t = 0$, $f'(0) = 1$) on obtient
\[
\log\frac{1-r}{(1+r)^3} \leq \log\abs{f'(z)} \leq
\log\frac{1+r}{(1-r)^3} :
\]
le théorème de distorsion, après exponentiation.

\textbf{14.} Supérieure~: le long du segment $[0, z]$,
\[
\abs{f(z)} = \Bigl|\int_0^rf'(t\eu^{\iu\theta})
\eu^{\iu\theta}\dd t\Bigr|
\leq \int_0^r\frac{1+t}{(1-t)^3}\dd t
= \frac{r}{(1-r)^2}
\]
($\frac{\dd}{\dd t}\frac t{(1-t)^2} = \frac{1+t}{(1-t)^3}$).
Inférieure~: noter $\frac r{(1+r)^2} < \frac14$ pour $r < 1$
(cela dit $(1-r)^2 > 0$), donc si $\abs{f(z)} \geq \frac14$
il n'y a rien à prouver. Sinon le segment $[0, f(z)]$ est
dans $D(0, \frac14) \subseteq f(\mathbb D)$ (question~7), et
$\gamma = f^{-1}\circ[0, f(z)]$ est un chemin $\mathcal C^1$
de $0$ à $z$ dans $\mathbb D$ ($f^{-1}$ \omterm{def:b3:holomorphic:holo}{holomorphe},
le~\cref{cor:b3:residues:openmapping}). En substituant $w =
f(\zeta)$,
\[
\abs{f(z)} = \int_{[0,f(z)]}\abs{\dd w}
= \int_\gamma\abs{f'(\zeta)}\,\abs{\dd\zeta}
\geq \int_0^1\frac{1 - \rho(s)}{(1 + \rho(s))^3}
\,\abs{\gamma'(s)}\,\dd s
\]
avec $\rho = \abs\gamma$, en utilisant la borne inférieure
de distorsion au rayon $\rho(s)$. Comme $\abs{\gamma'} \geq
\rho'$ (là où c'est défini~; $\rho$ est lipschitzienne) et
le facteur d'intégrande est positif,
\[
\abs{f(z)} \geq \int_0^1
\frac{1 - \rho(s)}{(1 + \rho(s))^3}\,\rho'(s)\,\dd s
= \int_0^{r}\frac{1 - u}{(1 + u)^3}\,\dd u
= \frac{r}{(1 + r)^2}
\]
($\rho(0) = 0$, $\rho(1) = r$~; la substitution n'utilise
qu'une primitive, $\frac{\dd}{\dd u}\frac u{(1+u)^2} =
\frac{1-u}{(1+u)^3}$, non la monotonie).

\textbf{15.} $k'(z) = \frac{\dd}{\dd z}\,z(1 - z)^{-2} =
\frac{1 + z}{(1 - z)^3}$. En $z = r$~: $k'(r) =
\frac{1+r}{(1-r)^3}$ et $k(r) = \frac r{(1-r)^2}$ --- les
deux bornes supérieures atteintes. En $z = -r$~: $k'(-r) =
\frac{1-r}{(1+r)^3}$ et $\abs{k(-r)} = \frac r{(1+r)^2}$
--- les deux bornes inférieures atteintes. La fonction de
Koebe étire son axe positif maximalement vers le bord
lointain et comprime le rayon antipodal maximalement vers
la pointe de sa fente.

\textbf{16.} $\abs{a_2} = 2$ signifie $\abs{b_1} = 1$ pour
$g(w) = 1/h(1/w) = w - \frac{a_2}2w^{-1} + \cdots$
(question~4)~; le théorème de l'aire $\sum n\abs{b_n}^2
\leq 1$ tue alors tout autre coefficient~: $g(w) = w + b_0
+ \eu^{\iu\alpha}/w$ avec $\eu^{\iu\alpha} = -\frac{a_2}2$.
Comme $h$ est impaire, $g(-w) = 1/h(-1/w) = -1/h(1/w) =
-g(w)$~: $g$ est impaire, donc $b_0 = 0$. En invertissant,
$h(z) = 1/g(1/z) = \frac{z}{1 + \eu^{\iu\alpha}z^2}$, et
$f(z^2) = h(z)^2$ donne
\[
f(w) = \frac{w}{(1 + \eu^{\iu\alpha}w)^2}
= \frac{w}{(1 - \eu^{\iu\theta}w)^2}
= \eu^{-\iu\theta}k\bigl(\eu^{\iu\theta}w\bigr),
\qquad \eu^{\iu\theta} = -\eu^{\iu\alpha} .
\]
Réciproquement chaque rotation a $a_2 = 2\eu^{\iu\theta}$ de
\omterm{def:b3:modules:module}{module} $2$~: les extrémales de Bieberbach sont exactement
les fonctions de Koebe tournées.

\textbf{17.} Si $c$ est omise avec $\abs c = \frac14$~: la
question~6 donne $F \in \mathcal S$ avec $A_2 = a_2 +
\frac1c$, donc $\frac1c = A_2 - a_2$ et
\[
4 = \Bigl|\frac1c\Bigr| = \abs{A_2 - a_2}
\leq \abs{A_2} + \abs{a_2} \leq 2 + 2 = 4 :
\]
égalité partout, en particulier $\abs{a_2} = 2$. Par la
question~16, $f = \eu^{-\iu\theta}k(\eu^{\iu\theta}\cdot)$,
dont l'ensemble omis est $\eu^{-\iu\theta}
\intoc{-\infty}{-\frac14}$~: l'unique valeur omise de
\omterm{def:b3:modules:module}{module} $\frac14$ est $c = -\eu^{-\iu\theta}/4$. (Contrôle~:
$a_2 = 2\eu^{\iu\theta}$ et $\frac1c = -4\eu^{\iu\theta}$,
donc $A_2 = -2\eu^{\iu\theta} = -a_2$~: l'inégalité
triangulaire est saturée par anti-alignement, comme il se
doit.)

\textbf{18.} Les \omterm{thm:b3:holomorphic:analytic}{estimations de Cauchy} sur le cercle $\abs z
= r$ (le~\cref{thm:b3:holomorphic:analytic}), combinées au
théorème de croissance, donnent
\[
\abs{a_n} \leq r^{-n}\sup_{\abs z = r}\abs f
\leq r^{1-n}\,(1 - r)^{-2} .
\]
Choisir $r = 1 - \frac1n$ ($n \geq 2$)~: $r^{1-n} =
\bigl(1 + \frac1{n-1}\bigr)^{n-1} < \eu$ (suite croissante
de limite $\eu$) et $(1 - r)^{-2} = n^2$~: $\abs{a_n} <
\eu\,n^2$.

\textbf{19.} $U = f(D(0, r))$ est \omterm{def:b3:topology:topology}{ouvert}
(le~\cref{cor:b3:residues:openmapping}) et contient $0$.
Bord~: si $p \in \partial U$, écrire $p = \lim f(z_k)$ avec
$z_k \in D(0, r)$~; une sous-suite donne $z_k \to z_\infty
\in \bar D(0, r)$, et $f(z_\infty) = p \notin U$ force
$z_\infty \in \partial D(0, r)$~: $\partial U \subseteq
f(\partial D(0, r))$, donc tout point du bord de $U$ a un
\omterm{def:b3:modules:module}{module} $\geq \frac r{(1+r)^2}$ (question~14). Soit
maintenant $\abs w < \frac r{(1+r)^2}$ et supposons $w
\notin U$. Le segment $[0, w]$ est \omterm{def:b3:topology:connected}{connexe}, rencontre $U$
(en $0$) et son complémentaire (en $w$), donc rencontre
$\partial U$~; mais tous ses points ont un \omterm{def:b3:modules:module}{module} $\leq \abs
w < \frac r{(1+r)^2}$~: contradiction. D'où $D\bigl(0,
\frac r{(1+r)^2}\bigr) \subseteq U$. Netteté~: $k(-r) =
-\frac r{(1+r)^2}$ est l'image d'un point du bord de
$D(0,r)$, et $k$ est injective, donc $k(-r) \notin k(D(0,
r))$~: le rayon ne peut être augmenté. Comme $r \to 1^-$,
$\frac r{(1+r)^2} \to \frac14$~: le théorème du quart pour
le disque entier.

\textbf{20.} D'abord, $\partial f(\mathbb D) \neq
\varnothing$~: sinon $f(\mathbb D)$, \omterm{def:b3:topology:topology}{ouvert}, fermé, non
vide, serait tout $\C$, et $f^{-1} \colon \C \to \mathbb D$
serait une fonction entière bornée non constante, contre
le~\cref{cor:b3:holomorphic:liouville}. Inégalité de
gauche~: la transformée de Koebe $F$ de la question~10 est
dans $\mathcal S$ et $F(\mathbb D) = \frac{f(\mathbb D) -
f(z_0)}{f'(z_0)(1-\abs{z_0}^2)}$~; par la question~7 elle
contient $D(0, \frac14)$, donc
\[
f(\mathbb D) \supseteq f(z_0) + D\Bigl(0,\
\tfrac14\abs{f'(z_0)}\bigl(1 - \abs{z_0}^2\bigr)\Bigr) ,
\]
et tout point de $\partial f(\mathbb D)$ --- disjoint de
l'\omterm{def:b3:topology:topology}{ouvert} $f(\mathbb D)$ --- est à distance $\geq
\frac14(1-\abs{z_0}^2)\abs{f'(z_0)}$ de $f(z_0)$. Inégalité
de droite~: soit $d = d(f(z_0), \partial f(\mathbb D)) <
\infty$. Le disque $D(f(z_0), d)$ est dans $f(\mathbb D)$~:
un segment de $f(z_0)$ à l'un de ses points reste à
distance $< d$ de $f(z_0)$, donc ne rencontre jamais
$\partial f(\mathbb D)$, et l'argument de connexité de la
question~19 le garde dans $f(\mathbb D)$. Alors $\hat g(w)
= f^{-1}(f(z_0) + dw)$ envoie $\mathbb D$ dans $\mathbb D$
avec $\hat g(0) = z_0$, et $\chi = \psi^{-1}\circ\hat g$,
avec $\psi(z) = \frac{z + z_0}{1 + \bar z_0z}$, fixe $0$~:
Schwarz (le~\cref{thm:b3:conformal:schwarz}) donne
$\abs{\chi'(0)} \leq 1$. Comme $\chi'(0) =
\frac{\hat g'(0)}{\psi'(0)} =
\frac{d}{f'(z_0)\,(1 - \abs{z_0}^2)}$, cela est $d \leq (1
- \abs{z_0}^2)\abs{f'(z_0)}$.

\textbf{21.} $k(\mathbb D) = \C\setminus
\intoc{-\infty}{-\frac14}$, donc $\partial k(\mathbb D) =
\intoc{-\infty}{-\frac14}$, et pour le point réel positif
$k(r) = \frac r{(1-r)^2}$ le point du bord le plus proche
est $-\frac14$~:
\[
d = k(r) + \frac14 = \frac{4r + (1-r)^2}{4(1-r)^2}
= \frac{(1+r)^2}{4(1-r)^2} .
\]
Membre de gauche de la question~20~: $\frac14(1 -
r^2)k'(r) = \frac14\,\frac{(1-r^2)(1+r)}{(1-r)^3} =
\frac{(1+r)^2}{4(1-r)^2} = d$~: égalité. (Le membre de
droite égale $4d$~: le facteur $4$ sépare les deux côtés, et
la fonction de Koebe siège exactement au bas.)

\textbf{22.} Le principe unique est \omterm{thm:b3:hilbert:parseval}{Parseval}~: l'univalence
interdit le chevauchement, donc l'aire du complémentaire de
l'image de $\{\abs w > \rho\}$, développée en modes de
Fourier sur les cercles, est non négative --- le théorème
de l'aire est une identité $L^2$ avec un signe. Tout le
reste est cette inégalité transportée~: une racine carrée
(question~3) la transforme en $\abs{a_2} \leq 2$~; une
réflexion de Möbius hors d'une valeur omise (question~6)
transforme $\abs{a_2} \leq 2$ en le théorème du quart~; les
automorphismes du disque (question~10) étendent $\abs{a_2}
\leq 2$ sur tout le disque comme théorème de distorsion~;
l'intégration radiale convertit la distorsion en
croissance, et la croissance en recouvrement. À chaque
étape le cas d'égalité se propage aussi, atterrissant
toujours sur les fonctions de Koebe tournées --- une
famille extrémale pour toute la théorie, tout comme les
rotations sont les uniques extrémales du lemme de Schwarz.
Une inégalité intérieure, plus la rigidité de son cas
d'égalité, gouverne toute la géométrie~: la
\omterm{def:b3:representations:rep}{représentation} conforme est l'art d'exploiter de telles
inégalités.

\textbf{23.} Le théorème de croissance borne $\abs f \leq
\frac{r}{(1-r)^2}$ uniformément sur chaque $\bar D(0, r)$,
$r < 1$, pour toutes les $f \in \mathcal S$ à la fois~:
localement bornée, donc $\mathcal S$ est une famille
normale (le~\cref{thm:b3:conformal:montel}). Si $f_n \in
\mathcal S \to f$ localement uniformément~: $f$ est
\omterm{def:b3:holomorphic:holo}{holomorphe} avec $f_n' \to f'$ localement uniformément
(Weierstrass), donc $f(0) = 0$, $f'(0) = 1$ --- en
particulier $f$ non constante --- et Hurwitz
(l'\cref{exo:b3:residues:8}) rend la limite d'applications
injectives injective~: $f \in \mathcal S$. Une
fonctionnelle \omterm{def:b3:topology:continuity}{continue} (telle que $f \mapsto \abs{a_2} =
\frac{\abs{f''(0)}}2$) sur une classe compacte
\emph{atteint} son supremum~: des fonctions extrémales
existent avant qu'on sache ce qu'elles sont --- le point de
départ de toute attaque variationnelle sur les problèmes de
coefficients, y compris Bieberbach.

\textbf{24.} Écrire $g(w) = w + A_2w^2 + O(w^3)$ et
composer~:
\[
z = g(f(z)) = f(z) + A_2f(z)^2 + O(z^3)
= z + (a_2 + A_2)z^2 + O(z^3),
\]
donc $A_2 = -a_2$ et $\abs{A_2} \leq 2$ (Bieberbach), avec
égalité ssi $\abs{a_2} = 2$, c'est-à-dire ssi $f$ est une
fonction de Koebe tournée (question~16) --- et alors $g$
est l'inverse correspondant, définie sur le plan fendu.

\textbf{25.} $f(z_1) - f(z_2) = (z_1 - z_2)\bigl(1 + a(z_1 +
z_2)\bigr)$. Si $\abs a \leq \frac12$~: pour $z_1 \neq z_2$
dans $\mathbb D$, $\abs{a(z_1 + z_2)} < 2\abs a \leq 1$
(strict~: $\abs{z_1 + z_2} < 2$), donc le second facteur ne
peut s'annuler~: injective, et $f \in \mathcal S$ (les
normalisations sont intégrées). Si $\abs a > \frac12$~: le
point $s = -\frac1a$ a $\abs s < 2$, donc $z_{1,2} =
\frac s2 \pm \varepsilon$ sont dans $\mathbb D$ pour
$\varepsilon > 0$ petit, sont distincts, et $z_1 + z_2 = s$
tue le facteur~: $f(z_1) = f(z_2)$, non injective. Donc
$\mathcal S$ contient $z + az^2$ exactement pour $\abs a
\leq \frac12$. La borne de Bieberbach $\abs{a_2} \leq 2$
est ainsi sauvagement non saturée par les polynômes de
\omterm{def:b3:galois:extension}{degré} $2$ --- les coefficients $a_n = n$ de la fonction de
Koebe viennent d'une série infinie conspirant le long du
rayon omis, un comportement qu'aucun polynôme (qui
n'appartient à $\mathcal S$ qu'avec de minuscules
coefficients) ne peut imiter.
\end{solution}
